\section*{Introduction to Inferential Statistics}

\emph{Inferential statistics} is the branch of statistics concerned with drawing conclusions about a population based on the analysis of a sample. Unlike descriptive statistics, which is limited to summarizing and describing observed data, statistical inference allows us to generalize and make predictions beyond the available data.

\subsection*{Why is Statistical Inference Necessary?}

In practice, we can rarely study an entire population of interest due to:

\begin{itemize}
	\item \textbf{Cost:} Studying every individual in a population can be prohibitively expensive.
	\item \textbf{Time:} Collecting data from the entire population can take too much time.
	\item \textbf{Accessibility:} Some members of the population may be inaccessible.
	\item \textbf{Destructive Testing:} In some cases, taking a measurement destroys the object being studied (for example, testing the lifespan of a lightbulb).
\end{itemize}

For these reasons, we take a \emph{representative sample} and use statistical techniques to infer the characteristics of the entire population.

\subsection*{Fundamental Concepts}

Before delving into inference techniques, it is crucial to understand the following concepts:

\begin{definicion}[Population and Sample]
	\begin{itemize}
		\item \textbf{Population:} The complete set of individuals, objects, or events sharing at least one common characteristic that we wish to study.
		\item \textbf{Sample:} A subset of the population selected for the study.
	\end{itemize}
\end{definicion}

\begin{definicion}[Parameter and Statistic]
	\begin{itemize}
		\item \textbf{Parameter:} A numerical measure that describes a characteristic of the \emph{population}. For example, the population mean $\mu$ or the population standard deviation $\sigma$.
		\item \textbf{Statistic:} A numerical measure calculated from a \emph{sample}. For example, the sample mean $\bar{x}$ or the sample standard deviation $s$.
	\end{itemize}
\end{definicion}

\begin{observacion}
	Parameters are fixed but generally unknown values. Statistics are random variables that vary from sample to sample and are used to estimate parameters.
\end{observacion}

\subsection*{Types of Statistical Inference}

There are two main types of statistical inference:

\begin{enumerate}
	\item \textbf{Estimation:} Consists of using sample statistics to estimate the value of a population parameter. It can be:
	\begin{itemize}
		\item \emph{Point estimation:} Providing a single value as the estimate of the parameter.
		\item \emph{Interval estimation:} Providing a range of values (a confidence interval) within which we expect the parameter to lie with a certain level of confidence.
	\end{itemize}
	
	\item \textbf{Hypothesis Testing:} Consists of formulating a claim (hypothesis) about a population parameter and using sample data to decide whether the evidence supports or contradicts that claim.
\end{enumerate}

\subsection*{The Statistical Inference Process}

The general process of statistical inference follows these steps:

\begin{enumerate}
	\item \textbf{Define the population} of interest and the characteristic to be studied.
	\item \textbf{Select a representative sample} using appropriate sampling techniques.
	\item \textbf{Calculate sample statistics} to serve as the basis for inference.
	\item \textbf{Apply inference techniques} (estimation or hypothesis testing) using appropriate probability distributions.
	\item \textbf{Interpret the results} in the context of the original problem, accounting for the uncertainty inherent in the sampling process.
\end{enumerate}

\subsection*{Sources of Error in Inference}

When drawing inferences from samples, we must consider two types of error:

\begin{itemize}
	\item \textbf{Sampling error:} The natural difference between a sample statistic and a population parameter due to the random variability inherent in the sampling process. This error decreases as the sample size increases.
	
	\item \textbf{Non-sampling error:} Systematic errors not caused by sampling, such as selection bias, measurement errors, non-response bias, etc. These errors do not decrease as the sample size increases and can be far more serious than sampling errors.
\end{itemize}

\begin{observacion}
	Inferential statistics is primarily concerned with quantifying and controlling sampling error. Controlling non-sampling errors requires careful study design and appropriate data collection procedures.
\end{observacion}

\subsection*{Sampling Distributions}

A fundamental concept in statistical inference is that of a \emph{sampling distribution}. If we take many samples of the same size from a population and calculate a statistic (such as the mean) for each sample, the probability distribution of these statistics is called the sampling distribution.

The \textbf{Central Limit Theorem} states that, under certain conditions, the sampling distribution of the sample mean approaches a normal distribution, regardless of the shape of the original population distribution, provided the sample size is sufficiently large.

This theorem is the cornerstone of many statistical inference techniques that we will study in this chapter.

\subsection*{Structure of the Chapter}

In the following sections we will study:

\begin{enumerate}
	\item \textbf{Random Sampling and the Central Limit Theorem:} Theoretical foundations for inference.
	\item \textbf{Hypothesis Testing:} Conceptual framework and methodology for making data-driven decisions.
	\item \textbf{Z and t Statistics:} Specific tools for hypothesis testing regarding means.
	\item \textbf{Confidence Intervals:} Method for estimating parameters with a quantified level of confidence.
	\item \textbf{Chi-Square Test:} Technique for analyzing categorical data and testing independence.
	\item \textbf{Correlation:} Measure of the linear relationship between variables.
\end{enumerate}

Throughout the chapter, we will use practical examples and Python code to illustrate the application of these concepts to real-world situations.
